\section{实施手册：把训练配方变成团队可执行流程}

\subsection{四周迭代节奏示例}
为了让本讲方法真正落地，最实用的做法是把训练配方变成固定节奏，而不是临时实验。下面给出一个四周节奏模板，可用于中等规模研究团队：
\begin{enumerate}
    \item 第 1 周：冻结目标任务与评估协议，完成数据去污染与基线跑通；
    \item 第 2 周：SFT 数据混合实验，确认行为稳定区间；
    \item 第 3 周：DPO/PPO 小规模扫描，验证偏好增益和副作用；
    \item 第 4 周：RLVR + test-time scaling 联合实验，输出部署成本曲线。
\end{enumerate}

\begin{importantbox}{让 ``实验节奏'' 比 ``灵感'' 更重要}
推理能力提升往往不是一次跳变，而是连续小改动的累积。将每周目标、评估指标和回归测试固定化，才能避免团队反复在同一坑里循环。
\end{importantbox}

\subsection{实验记录模板：最小可追溯字段}
\begin{table}[H]
\centering
\renewcommand{\arraystretch}{1.2}
\begin{tabular}{p{3.0cm}p{3.7cm}p{4.1cm}}
\toprule
\textbf{字段} & \textbf{记录内容} & \textbf{缺失后果} \\
\midrule
数据版本 & 样本来源、去重规则、配比、许可状态 & 无法判断增益是算法还是数据漂移 \\
训练配置 & 优化器、学习率、batch、训练步数、seed & 结果不可复现，难以回归 \\
奖励定义 & 偏好规则/验证器版本/阈值 & 难以排查 reward hacking \\
评估协议 & 基准集合、去污染规则、置信区间 & 分数不可比较，决策失真 \\
部署预算 & 平均延迟、P95、token 成本、失败重试率 & 无法判断是否可上线 \\
\bottomrule
\end{tabular}
\caption{将训练配方制度化的最小记录字段}
\end{table}

\begin{knowledgebox}{工程上最容易被忽视的一点}
很多团队只记录 ``最好成绩''，不记录 ``失败实验''。但在 reasoning 训练中，失败轨迹往往比成功轨迹更能揭示数据污染、奖励偏置或采样策略问题。
\end{knowledgebox}

\subsection{从研究到上线：灰度发布检查单}
\begin{warningbox}{上线前必须完成的四项检查}
\begin{itemize}
    \item \textbf{正确性灰度}：在高价值任务上先限制流量，验证真实通过率；
    \item \textbf{成本灰度}：对比单次解码与 test-time scaling 的收益/成本；
    \item \textbf{稳定性灰度}：观察长链输出是否出现格式漂移与重复推理；
    \item \textbf{安全灰度}：对高风险提示词与越权请求做红队回归。
\end{itemize}
\end{warningbox}

\subsection{本章小结}
训练配方能否产生长期价值，取决于是否被制度化为团队流程。固定节奏、完整记录和灰度检查，是把研究成果转成稳定产品能力的关键。


